\documentclass[12pt]{article}

\usepackage[a4paper, margin=2cm]{geometry}
\usepackage{fontspec}
\usepackage{polyglossia}
\usepackage{algorithm2e}
\usepackage{setspace}
\usepackage{tabularx, booktabs}
\usepackage{graphicx}
\usepackage{amsmath}
\usepackage{minted}
\usepackage{float}

\setmainfont{NewCM10-Regular}[
  Path,
  BoldFont        = NewCM10-Bold,
  ItalicFont      = NewCM10-Italic,
  BoldItalicFont  = NewCM10-BoldItalic,
]

\setsansfont{NewCMSans10-Regular}[
  Path,
  BoldFont        = NewCMSans10-Bold,
]

\setmonofont{NewCMMono10-Regular}[
  Path,
  BoldFont        = NewCMMono10-Bold,
  ItalicFont      = NewCMMono10-Italic,
]

\setdefaultlanguage{russian}

\newcolumntype{Y}{>{\centering\arraybackslash}X}

\begin{document}
\begin{titlepage}
\centering

{Министерство науки и высшего образования Российской Федерации}
\vskip 2mm
{\large\scshape Новосибирский государственный технический университет}
\vskip 2mm
{Кафедра прикладной математики}

\vfill
\vskip 3cm

{\scshape Лабораторная работа \textnumero 2}
\vskip 3mm
\begin{spacing}{1.7}
    {\huge\bf Метод конечных элементов для нелинейных уравнений}
\end{spacing}
\vskip 5mm
{\scriptsize ФПМИ, ПМ-24}
\vskip 2mm
{\large\it Параскун И., Герасименко В.}

\vfill

{преподаватели}
\vskip 2mm
{\scshape Вагин Д. В., д.т.н., доцент}
\vskip 1mm
{\scshape Задорожный А. Г., к.т.н., доцент}
\vskip 1mm
{\scshape Леонович Д. А., ассистент}

\vfill

{Новосибирск, \the\year}

\end{titlepage}

\setcounter{page}{2}

\section{Постановка задачи}

\noindent Построить МКЭ для нелинейного уравнения эллиптического типа

\begin{align}
  \label{equation} -\nabla(\lambda(u)\nabla u) + \gamma (u) u = f(u)
\end{align}

\vskip 4mm
\noindent с учётом следующих краевых условий:

\begin{align}
  &u \big|_{S_1} = u_g \\
  &\lambda(u)\frac{\partial u}{\partial n} \big|_{S_2} = \theta(u) \\
  &\lambda(u)\frac{\partial u}{\partial n} \big|_{S_3} +
    \beta(u \big|_{S_3} - u_{\beta}(u)) = 0
\end{align}

\section{Теоретическая часть}

\noindent Эквивалентная вариационная постановка в форме уравнения Галёркина для
уравнения \ref{equation}:

\begin{align}
  \begin{split}
    &\int_{\Omega}\lambda(u)\nabla u \nabla v_0 d\Omega +
    \int_{\Omega}\gamma(u)uv_0 d\Omega +
    \int_{S_3}\beta(u)uv_0 dS \\
    &= \int_{\Omega}f(u)v_0 d\Omega +
    \int_{S_2}\theta(u)v_0 dS +
    \int_{S_3}\beta(u)u_{\beta}(u)v_0 dS, \ \ \forall v_0 \in H_0^1
  \end{split}
\end{align}

\vskip 5mm
\noindent Раскладывая функции $u$ и $v_0$ по базису, переходим 
к конечноэлементной СЛАУ:

\begin{align}
  \begin{split}
    &\sum_{j=1}^{n}\big(
      \int_{\Omega}\lambda(u)\nabla \psi_j \nabla \psi_i d\Omega +
      \int_{\Omega}\gamma(u) \psi_j \psi_i d\Omega +
      \int_{S_3}\beta(u) \psi_j \psi_i dS
    \big)q_j \\
    &= \int_{\Omega}f(u) \psi_i d\Omega +
    \int_{S_2}\theta(u) \psi_i dS +
    \int_{S_3}\beta(u)u_{\beta}(u) \psi_i dS
  \end{split}
\end{align}

\vskip 5mm
\noindent или, в более удобной форме записи

\begin{align*}
  &A_{ij} = G_{ij} + M^{\gamma}_{ij} + M^{S_3}_{ij} \\
  &b_{i} = b^{\Omega}_{i} + b^{S_2}_{i} + b^{S_3}_{i}
\end{align*}

\vskip 5mm
\noindent где

\begin{align}
  &\hat{G}_{ij} = \int_{\hat{\Omega}}\hat{\lambda}(u)\nabla\hat{\psi}_j\nabla\hat{\psi}_i d\hat{\Omega} \\
  &\hat{M}^{\gamma}_{ij} = \int_{\hat{\Omega}}\hat{\gamma}(u)\hat{\psi}_j\hat{\psi}_i d\hat{\Omega} \\
  &\hat{M}^{S_3}_{ij} = \int_{\hat{S_3}}\hat{\beta}(u)\hat{\psi}_j\hat{\psi}_i d\hat{S} \\
  &\hat{b}^{\Omega}_{i} = \int_{\hat{\Omega}}\hat{f}(u)\hat{\psi}_i d\hat{\Omega} \\
  &\hat{b}^{S_2}_{i} = \int_{\hat{S_2}}\hat{\theta}(u)\hat{\psi}_i d\hat{S} \\
  &\hat{b}^{S_3}_{i} = \int_{\hat{S_3}}\hat{\beta}(u)\hat{u}_{\beta}(u)\hat{\psi}_i d\hat{S}
\end{align}

\vskip 5mm
\noindent Раскладывая параметры по базисным функциям на конечном элементе

\begin{align}
  &\hat{G}_{ij} = \int_{\hat{\Omega}}
    \sum_{k=1}^{8}(\hat{\lambda}_{k}\hat{\psi}_k)
    \nabla\hat{\psi}_j
    \nabla\hat{\psi}_i d\hat{\Omega} 
  = \sum_{k=1}^{8}\hat{\lambda}_{k}\int_{\hat{\Omega}}
    \hat{\psi}_k
    \nabla\hat{\psi}_j
    \nabla\hat{\psi}_i d\hat{\Omega} \\
  &\hat{M}^{\gamma}_{ij} = \int_{\hat{\Omega}}
    \sum_{k=1}^{8}(\hat{\gamma}_{k}\hat{\psi}_k)
    \hat{\psi}_j
    \hat{\psi}_i d\hat{\Omega} 
  = \sum_{k=1}^{8}\hat{\gamma}_{k}\int_{\hat{\Omega}}
    \hat{\psi}_k
    \hat{\psi}_j
    \hat{\psi}_i d\hat{\Omega} \\
  &\hat{M}^{S_3}_{ij} = \int_{\hat{S_3}}
    \sum_{k=1}^{4}(\hat{\beta}_{k}\hat{\psi}_k)
    \hat{\psi}_j
    \hat{\psi}_i d\hat{S}
  = \sum_{k=1}^{4}\hat{\beta}_{k}\int_{\hat{S_3}}
    \hat{\psi}_k
    \hat{\psi}_j
    \hat{\psi}_i d\hat{S} \\
  &\hat{b}^{\Omega}_{i} = \int_{\hat{\Omega}}
    \sum_{k=1}^{8}(\hat{f}_{k}\hat{\psi}_k)
    \hat{\psi}_i d\hat{\Omega} 
  = \sum_{k=1}^{8}\hat{f}_{k}\int_{\hat{\Omega}}
    \hat{\psi}_k
    \hat{\psi}_i d\hat{\Omega} \\
  &\hat{b}^{S_2}_{i} = \int_{\hat{S_2}}
    \sum_{k=1}^{4}(\hat{\theta}_{k}\hat{\psi}_k)
    \hat{\psi}_i d\hat{S} 
  = \sum_{k=1}^{4}\hat{\theta}_{k}\int_{\hat{S_2}}
    \hat{\psi}_k
    \hat{\psi}_i d\hat{S} \\
  &\hat{b}^{S_3}_{i} = \int_{\hat{S_3}}
    \sum_{k=1}^{4}(\hat{\beta}_{k}\hat{\psi}_k)
    \sum_{k=1}^{4}(\hat{u}_{\beta, k}\hat{\psi}_k)
    \hat{\psi}_i d\hat{S}
  = \sum_{k=1}^{4}\hat{\beta}_{k}\sum_{j=1}^{4}
    \hat{u}_{\beta, j}
    \int_{\hat{S_3}}\hat{\psi}_k\hat{\psi}_j\hat{\psi}_i d\hat{S}
\end{align}

\vskip 5mm
\noindent Учитывая вид трилинейных базисных функций

\begin{align*}
  & X_1(x) = \frac{x_{p+1}-x}{h_x}, & &X_2(x) = \frac{x - x_{p}}{h_x}, & &h_x = x_{p+1} - x_p, \\
  & Y_1(y) = \frac{y_{s+1}-y}{h_y}, & &Y_2(y) = \frac{y - y_{s}}{h_y}, & &h_y = y_{s+1} - y_s, \\
  & Z_1(z) = \frac{z_{r+1}-z}{h_z}, & &Z_2(z) = \frac{z - z_{r}}{h_z}, & &h_z = z_{r+1} - z_r.
\end{align*}

\begin{align}
  \psi_i = X_{\mu(i)}Y_{\nu(i)}Z_{\vartheta(i)}
\end{align}

\vskip 2mm
\noindent получаем следующие соотношения

\begin{align}
  \begin{split}
    \hat{G}_{ij} = \sum_{k=1}^{8}\hat{\lambda}_{k}(
      &{}^{n}G^{x}_{\mu(k)\mu(j)\mu(i)}{}^{n}M^{y}_{\nu(k)\nu(j)\nu(i)}{}^{n}M^{z}_{\vartheta(k)\vartheta(j)\vartheta(i)}+\\
      &{}^{n}M^{x}_{\mu(k)\mu(j)\mu(i)}{}^{n}G^{y}_{\nu(k)\nu(j)\nu(i)}{}^{n}M^{z}_{\vartheta(k)\vartheta(j)\vartheta(i)}+\\
      &{}^{n}M^{x}_{\mu(k)\mu(j)\mu(i)}{}^{n}M^{y}_{\nu(k)\nu(j)\nu(i)}{}^{n}G^{z}_{\vartheta(k)\vartheta(j)\vartheta(i)}
    )
  \end{split}
\end{align}

\begin{align}
  &\hat{M}^{\gamma}_{ij} = \sum_{k=1}^{8}\hat{\gamma}_{k}(
      {}^{n}M^{x}_{\mu(k)\mu(j)\mu(i)}{}^{n}M^{y}_{\nu(k)\nu(j)\nu(i)}{}^{n}M^{z}_{\vartheta(k)\vartheta(j)\vartheta(i)}
  ) \\
  &\hat{M}^{S_3}_{ij} = \sum_{k=1}^{4}\hat{\beta}_{k}(
      {}^{n}M^{\xi}_{\mu(k)\mu(j)\mu(i)}{}^{n}M^{\chi}_{\nu(k)\nu(j)\nu(i)}
  ) \\
  &\hat{b}^{\Omega}_{i} = \sum_{k=1}^{8}\hat{f}_{k}(
      M^{x}_{\mu(k)\mu(i)}M^{y}_{\nu(k)\nu(i)}M^{z}_{\vartheta(k)\vartheta(i)}
  ) \\
  &\hat{b}^{S_2}_{i} = \sum_{k=1}^{4}\hat{\theta}_{k}(
      M^{\xi}_{\mu(k)\mu(i)}M^{\chi}_{\nu(k)\nu(i)}
  ) \\
  &\hat{b}^{S_3}_{i} = \sum_{k=1}^{4}\hat{\beta}_{k}\sum_{j=1}^{4}\hat{u}_{\beta, j}(
      {}^{n}M^{\xi}_{\mu(k)\mu(j)\mu(i)}{}^{n}M^{\chi}_{\nu(k)\nu(j)\nu(i)}
  )
\end{align}

\vskip 5mm
\noindent где

\begin{align}
  &{}^{n}G^{x} = h_x
    \begin{pmatrix}
      \frac{1}{2} & -\frac{1}{2} \\
      -\frac{1}{2} & \frac{1}{2}
    \end{pmatrix} \\
  &{}^{n}M^{x} = \frac{1}{h_x}
    \begin{pmatrix}
      \frac{1}{4} & \frac{1}{12} \\
      \frac{1}{12} & \frac{1}{12}
    \end{pmatrix},
    \begin{pmatrix}
      \frac{1}{12} & \frac{1}{12} \\
      \frac{1}{12} & \frac{1}{4}
    \end{pmatrix}
\end{align}

\subsection{Метод простой итерации}

\noindent
Для решения полученной нелинейной системы организуем следующий итерационный процесс:

\begin{align}
  {\bf A}({\bf q}^{(k-1)}){\bf q}^{(k)} = {\bf b}({\bf q}^{(k-1)})
\end{align}

\vskip 5mm
\noindent
с условием остановки

\begin{align} \label{eq:1}
  R({\bf q}^{(k)}) = \frac{||{\bf A}({\bf q}^{(k)}){\bf q}^{(k)} - {\bf b}({\bf q}^{(k)})||}
  {||{\bf b}({\bf q}^{(0)})||} < \epsilon
\end{align}

\subsection{Метод Ньютона}

\noindent
Линеаризуем полученную нелинейную систему, раскладывая компоненты в ряд Тейлора:

\begin{align}
  &A_{ij}({\bf q})q_j \approx A_{ij}({\bf q}^{(k-1)})q_j^{(k-1)} +
    \sum_{r}\frac{\partial (A_{ij}({\bf q})q_j)}%
      {\partial q_r}\bigg\rvert_{{\bf q} = {\bf q}^{(k-1)}}(q_r - q_r^{(k-1)}) \\
  &b_i({\bf q}) \approx b_i({\bf q}^{(k-1)}) +
    \sum_{r}\frac{\partial b_{i}({\bf q})}%
      {\partial q_r}\bigg\rvert_{{\bf q} = {\bf q}^{(k-1)}}(q_r - q_r^{(k-1)})
\end{align}

\vskip 5mm
\noindent
Полагая, что параметр $\lambda$ зависит от решения, а остальные параметры нет,
получаем следующие выражения для составления линеаризованной системы:

\begin{align}
  &\hat{A}^{L}_{ij} = \hat{A}_{ij}(\hat{{\bf q}}^{(k-1)}) +
    \sum_{r=1}^{8}\frac{\partial \hat{A}_{ir}(\hat{\bf q})}%
      {\partial \hat{q}_j}\bigg\rvert_{\hat{\bf q} = \hat{\bf q}^{(k-1)}}
        \hat{q}_{r}^{(k-1)} \\
  &\hat{b}^{L}_{i} = \hat{b}_{i}(\hat{{\bf q}}^{(k-1)}) +
    \sum_{j=1}^{8}\hat{q}^{(k-1)}_{j}
      \sum_{r=1}^{8}\frac{\partial \hat{A}_{ij}(\hat{\bf q})}%
        {\partial \hat{q}_r}\bigg\rvert_{\hat{\bf q} = \hat{\bf q}^{(k-1)}}
          \hat{q}_{r}^{(k-1)} \\
  &\frac{\partial \hat{A}_{ij}(\hat{\bf q})}{\partial \hat{q}_r}
    \bigg\rvert_{\hat{\bf q} = \hat{\bf q}^{(k-1)}} =
    \int_{\hat{\Omega}}\frac{\partial\lambda^{\Omega}(\hat{\bf q}, \hat{\bf v}_r)}
        {\partial\hat{q}_r}\bigg\rvert_{\hat{\bf q} = \hat{\bf q}^{(k-1)}}
      \hat{\psi}_r\nabla\hat{\psi}_j\nabla\hat{\psi}_i d\Omega
\end{align}

\vskip 5mm
\noindent
Итерационный процесс организуем аналогично тому, как это делалось в методе
простой итерации.

\subsection{Использование параметра релаксации}

\begin{align}
  {\bf q}^{(k)} = \omega^{(k)}\bar{\bf {q}}^{(k)} 
    + (1 - \omega^{(k)}){\bf q}^{(k-1)}
\end{align}

\vskip 5mm
\noindent где $\omega^{(k)}$ определяется при помощи минимизации
функционала невязки \ref{eq:1}.

\newpage

\section{Описание тестирования программы}

\subsubsection{Описание задачи}

\noindent Пусть дана конечноэлементная сетка следующего вида:

\begin{figure}[H]
\centering
\includegraphics[width=\linewidth]{pde/test/non/case/img/msh.png}
\end{figure}

\noindent Для заданного решения $u(x, y, z) = x + y + z$ построим
краевую задачу

\begin{align*}
  &f = -\frac{\partial\lambda}{\partial x}
      -\frac{\partial\lambda}{\partial x}
      -\frac{\partial\lambda}{\partial x} + 0.4(x + y + z) \\
  &\theta\big|_{S_1} = -\lambda \\
  &\theta\big|_{S_2} = \lambda \\
  &\theta\big|_{S_3} = \lambda \\
  &u_{\beta}\big|_{S_4} = (x + y + z) - \frac{\lambda}{10} \\
  &u_{\beta}\big|_{S_5} = (x + y + z) + \frac{\lambda}{10} \\
  &u_{\beta}\big|_{S_6} = (x + y + z) - \frac{\lambda}{10}
\end{align*}

\vskip 5mm
\noindent при различных зависимостях $\lambda$. В случае задачи 
теплопроводности, решение выглядит следующим образом:

\begin{figure}[H]
\centering
\includegraphics[width=\linewidth]{pde/test/non/case/img/res.png}
\end{figure}

\subsubsection{Результат выполнения в различных условиях}

\begin{table}[H]
  \centering

  \begin{tabularx}{\textwidth}{@{}l|cYY|cYY@{}}
    \toprule
    $\lambda$ & \multicolumn{3}{c}{МПИ без релаксации} & \multicolumn{3}{c}{МПИ с релаксацией} \\
    \cmidrule(lr){2-4}
    \cmidrule(lr){5-7}
    & $k$ & $\overline{|u - u^*|}$ & $R({\bf q}^{(k)})$ & $k$ & $\overline{|u - u^*|}$ & $R({\bf q}^{(k)})$ \\
    \midrule
    $10$ & 1 & 5.8125638e-12 & 4.0489735e-14 & 1 & 5.8125638e-12 & 4.0489735e-14\\
    $u$ & 9 & 2.0090454e-06 & 3.4397768e-08 & 1 & 5.8125638e-12 & 4.0489735e-14\\
    $u^2$ & 17 & 3.1087130e-02 & 5.2876226e-08 & 1 & 5.8125638e-12 & 4.0489735e-14\\
    $u^3$ & 5 & - & - & 1 & 5.8125638e-12 & 4.0489735e-14\\
    $sin(u)$ & 1000 & 4.1737475e+01 & 1.9141056e-02 & 1 & 5.8125638e-12 & 4.0489735e-14\\
    \bottomrule
  \end{tabularx}
\end{table}

\section{Текст программы}

\inputminted{c}{pde/include/numx/pde/msh.h}
\inputminted{c}{pde/include/numx/pde/sim.h}
\inputminted{c}{pde/include/numx/pde/cnd.h}
\inputminted{c}{pde/include/numx/pde/slv.h}
\inputminted{c}{pde/include/numx/pde/fem.h}

\end{document}
